% =====================================================
% VISÃO GERAL DA DISCIPLINA
% =====================================================
\section{Visão Geral da Unidade Curricular}

Esta unidade curricular tem como objetivo consolidar e evoluir projetos desenvolvidos anteriormente no curso, promovendo a integração multidisciplinar de conhecimentos técnicos, metodológicos e científicos.

Ao longo do semestre, os estudantes irão revisar criticamente seus projetos anteriores, realizar diagnóstico técnico, propor ajustes estruturais, planejar nova execução e implementar melhorias por meio de micro-experimentos controlados.

A disciplina será orientada por uma perspectiva prática, investigativa e profissional, incentivando organização, autonomia, pensamento crítico e responsabilidade técnica na condução de projetos de tecnologia.

\subsection*{Estrutura do Semestre}

O semestre está organizado em 5 etapas estruturantes, distribuídas ao longo dos encontros:

\begin{enumerate}
    \item \textbf{Diagnóstico e Recuperação do Projeto} \\
    Revisão dos resultados anteriores, organização de repositórios, recuperação de datasets e análise crítica das limitações técnicas.

    \item \textbf{Plano de Ajustes e Replanejamento} \\
    Redefinição de objetivos, elaboração de cronograma atualizado, definição de métricas e metodologia de execução.

    \item \textbf{Execução Técnica e Implementação} \\
    Desenvolvimento de melhorias, refatoração de código, ajustes estruturais e consolidação da solução.

    \item \textbf{Micro-Experimentos e Validação} \\
    Planejamento e execução de 2 a 3 experimentos aplicados ao projeto, com coleta e análise de resultados.

    \item \textbf{Consolidação Científica e Apresentação Final} \\
    Produção de relatório técnico em formato de artigo científico utilizando \LaTeX{}, incluindo fundamentação teórica, trabalhos relacionados e apresentação dos resultados.
\end{enumerate}


% =====================================================
% FIGURA – PIPELINE DO PROJETO INTEGRADOR II
% =====================================================

\begin{figure}[H]
\centering
\resizebox{\linewidth}{!}{%
\begin{tikzpicture}[
    node distance=1.4cm and 1.2cm,
    every node/.style={font=\small},
    etapa/.style={
        rectangle,
        rounded corners=6pt,
        draw=azulTitulo,
        thick,
        minimum width=3.2cm,
        minimum height=1.55cm,
        text width=3.2cm,
        align=center,
        fill=azulTitulo!5
    },
    seta/.style={
        -{Latex[length=2.2mm]},
        thick,
        draw=azulTitulo
    }
]

\node[etapa] (e1) {\textbf{1. Diagnóstico}\\Recuperação do projeto\\Análise crítica};

\node[etapa, right=of e1] (e2) {\textbf{2. Plano de Ajustes}\\Replanejamento\\Definição de métricas};

\node[etapa, right=of e2] (e3) {\textbf{3. Execução Técnica}\\Implementação\\Refatoração};

\node[etapa, right=of e3] (e4) {\textbf{4. Micro-Experimentos}\\Validação\\Análise de dados};

\node[etapa, right=of e4] (e5) {\textbf{5. Consolidação Científica}\\Artigo em \LaTeX{}\\Apresentação final};

\draw[seta] (e1) -- (e2);
\draw[seta] (e2) -- (e3);
\draw[seta] (e3) -- (e4);
\draw[seta] (e4) -- (e5);

\end{tikzpicture}%
}
\caption{Pipeline estruturante do Projeto Integrador II.}
\end{figure}



\subsection*{Competências Desenvolvidas}

Ao final da disciplina, o estudante deverá ser capaz de:

\begin{itemize}
    \item Realizar análise crítica de projetos tecnológicos;
    \item Planejar e reorganizar projetos com base em diagnóstico técnico;
    \item Definir objetivos, métricas e cronogramas de execução;
    \item Executar melhorias técnicas estruturadas;
    \item Planejar e aplicar micro-experimentos com coleta de dados;
    \item Produzir documentação técnica estruturada em formato científico;
    \item Apresentar resultados técnicos de forma clara e fundamentada;
    \item Integrar conhecimentos multidisciplinares na solução de problemas reais.
\end{itemize}

\subsection*{Metodologia}

A disciplina será conduzida por meio de aprendizagem baseada em projetos (Project-Based Learning), com foco na evolução e maturação de projetos previamente desenvolvidos.

Serão realizados ciclos estruturados de diagnóstico, planejamento, execução e validação experimental, acompanhados por checkpoints periódicos.

Os grupos deverão manter um caderno de notas do projeto para registro contínuo das decisões técnicas, dificuldades, tutoriais, ajustes e aprendizados.

A documentação final será elaborada em formato de artigo científico utilizando template em \LaTeX{}, incentivando a cultura de organização técnica e possível submissão a evento acadêmico.

\subsection*{Avaliação da Aprendizagem}

A avaliação será processual, contínua e cumulativa, considerando a evolução técnica do projeto, a qualidade da documentação e o desempenho individual do estudante.

A nota final será composta por cinco componentes:

\begin{itemize}
    \item \textbf{Diagnóstico do Projeto (15\%)} \\
    Organização dos materiais anteriores, análise crítica e definição clara do estágio atual do projeto.

    \item \textbf{Plano de Ajustes e Replanejamento (20\%)} \\
    Definição estruturada de objetivos, cronograma, métricas e metodologia de execução.

    \item \textbf{Execução Técnica (30\%)} \\
    Implementação funcional das melhorias propostas e evolução concreta do projeto.

    \item \textbf{Micro-Experimentos e Análise de Resultados (20\%)} \\
    Planejamento, execução e interpretação adequada dos experimentos realizados.

    \item \textbf{Artigo Final e Apresentação (15\%)} \\
    Produção do relatório técnico em formato científico e apresentação oral estruturada.
\end{itemize}

A avaliação prioriza a maturidade técnica do projeto, a coerência metodológica, a capacidade de análise crítica e a qualidade da documentação produzida.
